\chapter{Impact Physics: From Club Delivery to Ball Launch}
\label{ch:impact}

The collision connects a golfer's delivered club state to the ball's launch velocity and spin.
A forward model predicts those outcomes from specified incident states, geometry and contact properties; an inverse model estimates some of those inputs from observations.
Either can be used with optical, radar or combined measurements.
The sensing modality alone does not establish a product's internal algorithm or which reported quantities are independently measured.
This chapter separates exact geometry, conditional impact mechanics and empirical approximations so that each can support an appropriate inference.

\section{The D-Plane}
\label{sec:dplane}

Use a right-handed basis with $x$ forward along the target, $y$ upward and $z$ to the right when viewed toward the target.
Horizontal angles remain positive to the right and elevations positive upward, consistent with the right-handed-player signs in \cref{ch:params}.
A unit direction with azimuth $\varphi$ and elevation $\lambda$ is
\begin{equation}
\label{eq:impactdirection}
\vect{d}(\varphi,\lambda)=
\begin{pmatrix}\cos\lambda\cos\varphi\\\sin\lambda\\\cos\lambda\sin\varphi\end{pmatrix}.
\end{equation}
The local face normal is $\uvec{n}=\vect{d}(\varphi_f,\lambda_d)$.
For collision mechanics, define $\uvec{p}_c=\vect{d}(\varphi_c,\alpha_c)$ from the incident velocity of the \emph{contact point}, with $V_c=\|\vect{v}_c\|>0$ and all states evaluated at a declared common instant.
A monitor's reported club path and attack angle may instead describe another reference point or maximum compression.
They cannot be substituted silently: for a rigid head, $\vect{v}_c=\vect{v}_O+\boldsymbol\omega_h\times(\vect{r}_c-\vect{r}_O)$, and the state can change during impact.

In the centered, symmetric oblique-impact idealization associated with the D-plane~\cite{jorgensen}, the nonparallel directions $\uvec{n}$ and $\uvec{p}_c$ span a plane.
Assume an initially stationary, nonspinning spherical ball, an isotropic contact response and a face normal treated as fixed during the impulse.
If the net impulse lies in that plane, so does the outgoing ball velocity.
Its location between face and path, and its closeness to either, additionally depend on the contact law; geometry alone supplies no universal weighting.
For tangential impulse directed along the contact-plane component of club motion, the ball's contact lever arm is $-R\uvec{n}$, giving
\begin{equation}
\label{eq:spinaxisdplane}
\Delta\boldsymbol\omega_b
 = I_b^{-1}(-R\uvec{n})\times\vect{J}_t,
\qquad
\Delta\boldsymbol\omega_b\parallel\uvec{p}_c\times\uvec{n}.
\end{equation}
Here $R$ is the ball radius and $I_b$ its scalar inertia in this spherical model; the signed tangential impulse sets the orientation along that axis and its magnitude.
Parallel directions do not define a unique D-plane or a normalized cross-product axis.
Initial spin, off-center head rotation, anisotropy and changing contact geometry require an expanded model.

\begin{figure}[htbp]
\centering
\begin{tikzpicture}[scale=0.95,>=Stealth]
  \draw[accentgray,dashed,->] (0,0) -- (9.5,0);
  \node[accentgray] at (5.2,0.8) {\small Target Line};
  \draw[accentgray,thin] (5.2,0.55) -- (5.2,0.04);
  \draw[thick,accentblue,->] (0,0) -- ({8*cos(-2)},{8*sin(-2)});
  \draw[accentblue,thin] ({8*cos(-2)},{8*sin(-2)}) -- (8.4,0.45);
  \node[accentblue,anchor=west] at (8.4,0.6)
    {\small Face Normal ($-2\degs$)};
  \draw[very thick,red!70!black,->] (0,0) -- ({8.6*cos(-2.96)},{8.6*sin(-2.96)});
  \node[red!70!black,anchor=west] at (8.8,-0.55)
    {\small Launch ($-2.96\degs$)};
  \draw[thick,accentgreen,->] (0,0) -- ({8*cos(-6)},{8*sin(-6)});
  \draw[accentgreen,thin] ({8*cos(-6)},{8*sin(-6)}) -- (8.4,-1.25);
  \node[accentgreen,anchor=west] at (8.4,-1.4)
    {\small Club Path ($-6\degs$)};
\end{tikzpicture}
\caption{Manufactured Horizontal Launch Example.
The declared $0.76/0.24$ weighting gives $-2.96\degs$ from face angle $-2\degs$ and path $-6\degs$.
This top-view projection illustrates the arithmetic, not a measured shot or a complete three-dimensional spin prediction.}
\label{fig:dplane}
\end{figure}

\subsection{Face/Path Weighting of Launch Direction}
\label{sec:facepathweighting}

The filtered player data of Wood~et~al.~\cite{facepathstudy} give the
following study-specific horizontal weights in an approximate model of
launch direction between face angle and club path:
\begin{equation}
\label{eq:launchweight}
\varphi_{\mathrm{launch}} \;\approx\;
  w_f\,\varphi_f + (1-w_f)\,\varphi_p,
\qquad
w_f \approx
\begin{cases}
  0.76 \pm 0.08 & \text{driver} \\
  0.69 & \text{7-iron} \\
  0.61 & \text{wedge.}
\end{cases}
\end{equation}

\begin{warning}
The values above are the \emph{horizontal} weights measured by
Wood~et~al.~\cite{facepathstudy}. Their player data contain 1{,}575 shots
before filtering, with 157 golfers testing driver/7-iron and ten also
testing a wedge. Club motion was measured optically at 720\,fps and ball
launch with a GC2. The horizontal means use strikes within 0.25\,in of
face centre in each axis and at least $1.5\degs$ face-to-path separation:
89 driver, 172 iron and 28 wedge shots remained. A separate 15-shot PING~Man
robot test returned $0.63$ for the 7-iron.

The familiar $0.85/0.15$ figure comes from TrackMan's \emph{Ten
Fundamentals}, which gives driver rules of thumb in both planes (and
$75\%$ for irons). Wood's driver means are $0.83 \pm 0.08$ vertically
and $0.76 \pm 0.08$ horizontally, reported with 95\% confidence intervals.
The difference from 0.85 is a reason to retain the model's conditions,
not a universal calibration correction. These are estimates from filtered
samples, not individual-shot tolerance bands; protocol and measurand
differences preclude treating the comparison as a definitive test of every
application of the vendor rule.

For any estimator using this weighting model, inverting
Eq.~\eqref{eq:launchweight} to recover face angle amplifies any bias in
measured launch direction by $1/w_f$ when path and weight are fixed.
At $w_f=0.76$ the factor is $1.32$, versus $1.18$ at $w_f=0.85$:
about 12\% greater sensitivity, not a factor of two. Path error and
uncertainty in the weight contribute separately, so this derivative is
not a complete uncertainty budget for inferred face angle.
\end{warning}

Henrikson and colleagues~\cite{mdpifriction} compare a tangentially compliant contact model with plate impacts and player data.
Their model's launch ratio varies with incidence even at fixed friction, and changing friction has different effects across incidence regimes.
It uses a fixed impact surface and an adjustment for normal inelasticity; agreement in those comparisons is not a universal validation for moving, rotating clubheads.
The findings support examining the contact-force history rather than attributing all club-dependent weighting to surface smoothness.

\subsection{Spin Loft and Approximate Launch Maps}

The geometric angle $\Phi$ between the local normal and contact-point incident velocity satisfies the full dot product
\begin{equation}
\label{eq:obliqueness}
\cos\Phi=\uvec{n}\cdot\uvec{p}_c
 =\sin\lambda_d\sin\alpha_c
  +\cos\lambda_d\cos\alpha_c\cos(\varphi_f-\varphi_c).
\end{equation}
This contact obliquity equals a reported spin loft only when point, time and angle definitions match.
For coplanar forward directions it reduces to $|\lambda_d-\alpha_c|$.
Writing $A_c=\varphi_f-\varphi_c$, neither $\cos A_c\cos\lambda_d$ nor $\cos A_c\cos(\lambda_d-\alpha_c)$ is a general replacement when horizontal separation and both elevations are nonzero.

Tutelman's historical model~\cite{tutelman3d} uses its own loft/relative-angle construction and an approximate launch map,
\begin{equation}
\label{eq:tutelmanlaunch}
f(\Phi)=0.96-0.0071\Phi,\qquad
\mathrm{DA}_{\rm vert}=L f(\Phi),\qquad
\mathrm{DA}_{\rm horiz}=A f(\Phi),
\end{equation}
with angles in degrees and $A,L$ interpreted according to that model.
Its historical loft terminology is not interchangeable with target-frame dynamic loft, and its later comparison with TrackMan does not establish that the original model was fitted to that comparison.
Treat the map as an author approximation with a declared domain, not an exact spherical-coordinate transformation or a general friction identification.
For a three-dimensional implementation, obtain launch angles from the predicted velocity vector using \texttt{atan2}; adding component angles is not a general rotation composition.

\subsection{Spin-Axis Tilt}

A spin axis needs a reference frame as well as a direction.
For outgoing unit velocity $\uvec{v}_b$ that is not vertical, let $\uvec{e}_y$ be world up and define
\begin{equation}
\label{eq:axistilt}
\uvec{e}_b=\frac{\uvec{v}_b\times\uvec{e}_y}
 {\|\uvec{v}_b\times\uvec{e}_y\|},\quad
\uvec{e}_u=\uvec{e}_b\times\uvec{v}_b,\quad
\theta_{\rm axis}=\operatorname{atan2}
 (-\boldsymbol\omega_b\cdot\uvec{e}_u,
   \boldsymbol\omega_b\cdot\uvec{e}_b).
\end{equation}
This convention makes positive backspin point along $\uvec{e}_b$ and gives positive tilt for a right-curving spin component in ordinary forward flight.
Compare a device's sign and reference-axis definition before comparing its number.
A vertical launch needs a different frame construction; a zero spin vector has no spin-axis direction.
The tilt is also undefined when both transverse projections vanish, even if an axial spin component is nonzero.
A general spin vector can also have a component along $\uvec{v}_b$, so total spin and one tilt angle do not specify every three-dimensional spin state.

The often-used ratio $\tan(\text{face-to-path})/\tan(\text{spin loft})$ is not an exact tilt identity for general target-frame angles.
Calculate the cross product and its projections under the stated collision assumptions instead.
Likewise, a fixed yards-per-degree or return-to-target multiplier cannot replace a flight calculation: speed, spin magnitude, launch, wind and aerodynamic coefficients also affect curvature.

\begin{implication}
If an estimator adopts the launch-weight model, it may infer
$\varphi_f=[\varphi_{\rm launch}-(1-w_f)\varphi_p]/w_f$.
Its local launch sensitivity is $1/w_f$; path, weight, reference-frame covariance and model discrepancy contribute separately.
Agreement between an inferred face angle and the launch measurement used to infer it is not an independent validation of that face angle.
\end{implication}

\section{Ball Speed and Smash Factor}
\label{sec:smash}

Let $J_n$ be the normal impulse on an initially stationary ball of mass $m$.
For a free rigid head of mass $M$, centroidal inertia tensor $I_h$ and contact lever arm $\vect{r}_h$, its normal inverse effective mass is
\begin{equation}
\label{eq:impacteffectivemass}
\frac{1}{M_{\rm eff,n}}=\frac{1}{M}
 +(\vect{r}_h\times\uvec{n})^{\mathsf T}I_h^{-1}
   (\vect{r}_h\times\uvec{n}).
\end{equation}
In the scalar limit with no normal--tangential coupling, normal restitution $e_n$ and closing speed $V_c\cos\Phi>0$ give~\cite{cochranstobbs,penner2003}
\begin{equation}
\label{eq:ballspeed}
J_n=\frac{(1+e_n)V_c\cos\Phi}{1/m+1/M_{\rm eff,n}},
\qquad
\vect{v}_b^+\cdot\uvec{n}
 =\frac{1+e_n}{1+m/M_{\rm eff,n}}V_c\cos\Phi.
\end{equation}
This predicts the \emph{normal component} of ball velocity.
If there is tangential impulse, $\vect{v}_b^+=(J_n\uvec{n}+\vect{J}_t)/m$ and the speed norm also contains that component.
Normal--tangential coupling requires a coupled contact solve rather than inserting an effective mass into an otherwise unrelated scalar fit.

For illustration only, taking $m=0.04593$ kg, $M_{\rm eff,n}=0.200$ kg and $e_n=0.83$ gives the following normal-component ratios to $V_c$:
\begin{center}
\small
\begin{tabular}{@{}lcccc@{}}
\toprule
Obliquity & $0\degs$ & $10\degs$ & $20\degs$ & $30\degs$ \\
Normal Ratio & 1.488 & 1.466 & 1.398 & 1.289 \\
\bottomrule
\end{tabular}
\end{center}
These are manufactured values, not a current equipment-conformance test or a loft-dependent physical ceiling.
The assumed restitution is a model parameter; it does not identify every real ball/face combination or impact speed.

\begin{keypoint}
A reported smash factor divides measured ball speed by a specified reported club speed, which need not be the contact-point incident speed $V_c$.
Exceeding one row of this illustrative table therefore does not by itself prove measurement error.
Before rejecting a shot, reconcile reference point and time, effective inertia, contact model and uncertainty.
A discrepancy is a diagnostic lead, not a universal impossibility verdict.
\end{keypoint}

\section{Spin Generation}
\label{sec:spingen}

The same tangential impulse changes both ball translation and angular momentum.
For a spherical ball with contact arm $-R\uvec{n}$,
$\Delta\boldsymbol\omega_b=I_b^{-1}(-R\uvec{n})\times\vect{J}_t$.
For an impulse in the obliquity plane, the launch deflection from the normal satisfies $\tan\delta_b=J_t/J_n$.
Thus independently chosen launch and spin fits must still be checked for compatibility with one contact solution.
Tangential compliance can reverse force during contact; a single friction coefficient does not specify the entire slip--stick history~\cite{mdpifriction,crossoblique}.

When its club-speed input is identified with $V_c$, an empirical approximation in Tutelman's model is~\cite{tutelman3d}
\begin{equation}
\label{eq:spinrate}
S\ [\mathrm{rpm}]\approx160\,V_c[\mathrm{mph}]\sin\Phi.
\end{equation}
The coefficient carries the stated units and is not a dimensional identity or a universal spin law.
Retain the model's input conventions and calibration limits; use the predicted spin vector for component decomposition rather than imposing a separate trigonometric tilt rule.
Rules based only on iron loft, or a universal spin-loft saturation angle, omit contact material, speed and force history.
Changing surface conditions changes the contact problem; this chapter supplies no universal contamination correction or threshold.
% Retain the previously cited vendor source as background, not as evidence of a universal threshold.
\nocite{trackmanspinloft}

\section{Gear Effect and Impact Location}
\label{sec:gear}

An off-center impulse rotates the head and changes contact-point velocity during collision.
In a simplified horizontal construction, let $x$ be the signed lateral offset from the head CG line, $C$ the normal CG-to-face distance and $I_h$ the head inertia about the relevant vertical axis.
Tutelman's angular-impulse heuristic~\cite{tutelmangear} estimates $\Delta\omega_h=xmV_b/I_h$ and equates a face tangential speed to the ball's rotational surface speed at separation, giving
\begin{equation}
\label{eq:gear}
s\ [\mathrm{rpm}]\approx58{,}830\,
 \frac{V_b[\mathrm{mph}]\,C[\mathrm{in}]\,x[\mathrm{in}]}
 {I_h[\mathrm{g\,cm^2}]}
 \approx16.4\,V_b[\mathrm{mph}]\,x[\mathrm{in}].
\end{equation}
The last coefficient uses the author's representative $C=1.42$ in and $I_h=5100\,\mathrm{g\,cm^2}$; there is no extra factor of 100.
For $V_b=150$ mph and $x=0.25$ in, the simplified estimate is 615 rpm, not 61,500 rpm.
The cited 2.5\% describes the relative sample standard deviation of a historical, filtered inertia/depth ratio, not a guaranteed spin-prediction error or all current drivers.
The contact-speed assumption neglects ball tangential translation and coupled recoil, so the heuristic is not a complete no-slip impact solution.

\subsection{Why a Coupled Contact Solve Matters}

The missing terms can be seen without fitting a golfer or product.
Consider an initially nonspinning free head moving normally at speed $V$ toward a stationary spherical ball, with head contact arm $C\uvec{n}+x\uvec{t}$ in a planar model.
The unit tangent $\uvec{t}$ is perpendicular to $\uvec{n}$.
The head has mass $M$ and scalar inertia $I_h$; the ball has mass $m$, radius $R$ and inertia $I_b$.
Take positive inertial parameters and $V>0$, with $0\leq e_n\leq1$ and an impulsive friction coefficient $\mu\geq0$.
With impulse components $(J_n,J_t)$ on the ball, the ball-minus-head contact-velocity increment is $K(J_n,J_t)^{\mathsf T}$, where
\begin{equation}
\label{eq:gearcoupledmobility}
K=\begin{pmatrix}
 \frac1m+\frac1M+\frac{x^2}{I_h} & -\frac{Cx}{I_h}\\
 -\frac{Cx}{I_h} & \frac1m+\frac1M+\frac{R^2}{I_b}+\frac{C^2}{I_h}
\end{pmatrix}.
\end{equation}
Normal restitution and zero final tangential relative velocity then give the sticking candidate
\begin{equation}
\label{eq:gearstickcandidate}
K\begin{pmatrix}J_n\\J_t\end{pmatrix}
 =\begin{pmatrix}(1+e_n)V\\0\end{pmatrix},
\qquad |J_t|\leq\mu J_n.
\end{equation}
The inequality must hold for the candidate to be admissible under this impulsive friction model; otherwise a sliding branch is needed.
The term $1/m$ in the tangential entry accounts for ball translation, $R^2/I_b$ for ball spin and the head terms for recoil.
Consequently, matching contact velocities does not mean that face speed equals $R$ times ball spin.
This idealization still omits contact compliance, shaft impulses and changing face geometry; it is a consistency calculation, not a universally identified gear-effect law.

Bulge and roll alter the local normal as strike location changes, while gear effect alters the impulse through head motion.
Their contributions must be evaluated together; changing the rotation axis also changes the relevant head inertia.
An optical strike measurement can constrain the geometry, but uncertainty in pose, CG, inertia and contact response remains.
Neither a quoted yardage example nor \cref{eq:gear} provides a universal bound on face-angle inference error.

\begin{implication}
Impact location is valuable when it resolves an important uncertainty for the intended outcome.
Its priority relative to face pose, local velocity or spin requires a sensitivity and uncertainty analysis.
Off-center launch/spin cannot in general be inverted with a centered D-plane model alone, and adding a strike sensor does not remove collision-model discrepancy.
\end{implication}

\section{Vertical Plane and the Low-Point Geometry}

A fixed inclined-plane construction explains one conditional coupling between attack and horizontal path.
Set the plane's horizontal base along the target and its inclination to $0<\beta<90\degs$, with unit normal $(0,\cos\beta,\sin\beta)$ in the forward/up/right frame.
A velocity direction with elevation $\alpha_c$ and horizontal offset $\psi$ lies in that plane only if
\begin{equation}
\label{eq:inclinedplanepath}
\sin\alpha_c\cos\beta+
 \cos\alpha_c\sin\psi\sin\beta=0,
\qquad
\sin\psi=-\frac{\tan\alpha_c}{\tan\beta}.
\end{equation}
On the forward-moving branch, $\alpha_c=-5\degs$ and $\beta=60\degs$ give $\psi\approx2.8953\degs$ to the right.
The result requires $|\tan\alpha_c/\tan\beta|\leq1$ and the declared side of plane inclination; a mirrored plane changes the sign.
A measured swing need not follow one fixed plane, and this derivation does not establish any particular product's processing algorithm.

\section{Connecting Delivery, Impact and Measurement}

Swing dynamics determines the delivered local velocity, orientation and retained elastic state; contact mechanics converts that state into an impulse and outgoing ball state; the instrument estimates selected quantities from its observations.
An explanation of technique must keep those stages connected while testing each with sufficiently independent evidence.
For example, changing wrist motion may alter both face pose and contact-point velocity, and changing strike location may alter both the local normal and effective inertia.
A single fitted launch percentage cannot identify which change caused an observed ball flight.
Report the retained states, reference points, contact assumptions and uncertainty alongside the conclusion, and compare joint predictions of launch direction, speed and spin rather than validating each with incompatible scalar rules.
